% BHU PYQ -- Manually transcribed & error-corrected
\documentclass[11pt,a4paper]{article}

\usepackage[a4paper,top=2.5cm,bottom=2.5cm,left=2.8cm,right=2.8cm,headheight=14pt]{geometry}
\usepackage{amsmath,amssymb}
\usepackage[T1]{fontenc}
\usepackage[utf8]{inputenc}
\usepackage{lmodern}
\usepackage{microtype}
\usepackage{fancyhdr}
\usepackage{enumitem}
\usepackage{hyperref}
\usepackage{lastpage}
\usepackage{parskip}

\hypersetup{colorlinks=true,urlcolor=black,linkcolor=black,
  pdftitle={CHB-504: Physical Chemistry-III -- BHU PYQ},pdfauthor={Banaras Hindu University}}

\pagestyle{fancy}\fancyhf{}
\renewcommand{\headrulewidth}{0.4pt}\renewcommand{\footrulewidth}{0.4pt}
\fancyhead[L]{\small\textit{Banaras Hindu University}}
\fancyhead[R]{\small\textit{CHB-504: Physical Chemistry-III}}
\fancyfoot[C]{\small Page \thepage\ of \pageref{LastPage}}

\newlist{parts}{enumerate}{2}
\setlist[parts,1]{label=(\alph*),leftmargin=*,itemsep=5pt,topsep=3pt}
\setlist[parts,2]{label=(\roman*),leftmargin=*,itemsep=3pt,topsep=2pt}
\newcommand{\pts}[1]{\hfill\textit{[#1]}}

\begin{document}

\begin{center}
  {\large\bfseries BANARAS HINDU UNIVERSITY}\\[2pt]
  {\normalsize B.Sc. (Hons.) Semester V Examination, 2016-17}\\[6pt]
  \rule{0.8\linewidth}{0.5pt}\\[6pt]
  {\large\bfseries Paper No. CHB-504: Physical Chemistry-III}\\[4pt]
  {\normalsize Time: 3 Hours\hfill Maximum Marks: 70}
\end{center}
\rule{\linewidth}{0.4pt}
\small
\noindent\textit{Instructions: Answer five questions in total, including Question~1 which is compulsory.}
\normalsize
\bigskip

\noindent\textbf{Question 1.} Answer all parts of the following: \pts{5 $\times$ 2 = 10}
\begin{parts}
  \item Write the BET equation and identify each term involved in it.
  \item What is meant by the ionic strength of a solution?
  \item Explain why there is a depression in freezing point of a solvent upon addition of a non-volatile solute.
  \item State the Grotthuss-Draper law and Stark-Einstein law of photochemistry.
  \item What is quantum yield? Write the reasons for high and low quantum yields.
\end{parts}

\medskip\rule{\linewidth}{0.2pt}

\noindent\textbf{Question 2.}
\begin{parts}
  \item What is the photostationary state? Deduce the Stern-Volmer equation and show how the value of the quenching constant and other constants could be determined. \pts{8}
  \item Derive the rate law for a unimolecular surface reaction on a solid catalyst (Langmuir-Hinshelwood mechanism). \pts{6}
\end{parts}

\medskip\rule{\linewidth}{0.2pt}

\noindent\textbf{Question 3.}
\begin{parts}
  \item Discuss the influence of high-frequency AC on the conductance of electrolyte solutions (Falkenhagen effect). \pts{6}
  \item Explain the reduction in ionic mobility due to 'relaxation' and 'electrophoretic' effects. Why do these effects become significant with increasing concentration? \pts{8}
\end{parts}

\medskip\rule{\linewidth}{0.2pt}

\noindent\textbf{Question 4.}
\begin{parts}
  \item Derive the Debye-H\"uckel-Onsager equation for the conductance of a strong electrolyte solution at a given concentration. Write the limiting forms. \pts{8}
  \item Discuss the primary and secondary processes in photochemical reactions. \pts{6}
\end{parts}

\medskip\rule{\linewidth}{0.2pt}

\noindent\textbf{Question 5.} Write short notes on any two of the following: \pts{2 $\times$ 7 = 14}
\begin{parts}
  \item Photosensitized reactions
  \item Adsorption isotherms (Langmuir vs. Freundlich)
  \item Wien effect on conductance under high electric fields
\end{parts}

\vfill
\rule{\linewidth}{0.4pt}
\begin{center}\small
  \href{https://bhu-exam.vercel.app/}{Click here to visit our website}\\[4pt]
  \href{https://me-aryan.vercel.app/}{Click here to visit our contributor}\\[4pt]
  \href{https://quashan-me.vercel.app/}{Click here to visit our contributor}
\end{center}

\end{document}
